


\section{Proofs and Additional Reformulation Details}
\label{app:sg_proofs}
\section{Proofs of the Results in Section~\ref{sec:solution_methodology_compact}}
\label{app:sg_proofs}

For notational brevity, write $\Pi:=\bar\Pi$ and recall that
\[
\Omega(K):=\left\{\delta\in\{0,1\}^{|\mathcal L|}:\ \mathbf 1^\top\delta\le K\right\}.
\]
We also use the shorthand
\[
V:=\mathrm{ext}(\Pi)
\]
for the set of extreme points of the dual polyhedron \eqref{eq:dual_poly_compact}.

\subsection{Proof of Proposition~\ref{prop:sg_moment_dual}}

\begin{proof}[Proof of Proposition~\ref{prop:sg_moment_dual}]
Because $\Omega(K)$ is finite, write
\[
\Omega(K)=\{\delta^1,\ldots,\delta^S\}.
\]
Any distribution $\mathbb P\in\mathfrak P$ can then be represented by a probability vector
$p\in\mathbb R_+^S$ satisfying $\sum_{s=1}^S p_s=1$ and
\[
\sum_{s=1}^S p_s\delta^s\le \mathbf{FP}.
\]
Therefore, for fixed $x$,
\[
\begin{aligned}
    &w_B(x)=\\
    &\max_{p\ge 0}
    \left\{
    \sum_{s=1}^S p_s f_{x,B}(\delta^s):
    \sum_{s=1}^S p_s=1,\
    \sum_{s=1}^S p_s\delta^s\le \mathbf{FP}
    \right\}.
\end{aligned}
\]
This is a finite-dimensional linear program. Let $\alpha\in\mathbb R$ be the dual variable associated with
$\sum_{s=1}^S p_s=1$, and let $\lambda\ge 0$ be the dual multiplier associated with the componentwise constraint
$\sum_{s=1}^S p_s\delta^s\le \mathbf{FP}$. By linear-programming strong duality,
\[
\begin{aligned}
    &w_B(x)=\\
    &\min_{\alpha\in\mathbb R,\ \lambda\ge 0}
    \left\{
    \alpha+\lambda^\top \mathbf{FP}:
    \alpha+\lambda^\top\delta^s\ge f_{x,B}(\delta^s), s=1,\ldots,S
    \right\},
\end{aligned}
\]
which is exactly \eqref{eq:sg_moment_dual}.

Next, substitute the representation \eqref{eq:sg_fx_dual} of $f_{x,B}(\delta)$ into
\eqref{eq:sg_moment_dual}. The inequality
\[
\alpha+\lambda^\top\delta\ge f_{x,B}(\delta)
\qquad \forall\delta\in\Omega(K)
\]
is equivalent to
\[
\alpha+\lambda^\top\delta\ge
\beta(\boldsymbol\pi)-\gamma(\boldsymbol\pi)^\top x+\phi(\boldsymbol\pi)^\top\delta,
\qquad \forall \delta\in\Omega(K),\ \forall\boldsymbol\pi\in\Pi^B.
\]
Rearranging yields
\[
\alpha\ge
\beta(\boldsymbol\pi)-\gamma(\boldsymbol\pi)^\top x+
\big(\phi(\boldsymbol\pi)-\lambda\big)^\top\delta,
\quad \forall \delta\in\Omega(K),\ \forall\boldsymbol\pi\in\Pi^B.
\]
Taking the maximum over $\delta\in\Omega(K)$ gives the equivalent form
\[
\alpha
\ge
\beta(\boldsymbol\pi)-\gamma(\boldsymbol\pi)^\top x
+\max_{\delta\in\Omega(K)}\big(\phi(\boldsymbol\pi)-\lambda\big)^\top\delta,
\quad \forall\boldsymbol\pi\in\Pi^B.
\]
The proof is complete.
\end{proof}

\subsection{Proof of Lemma~\ref{lem:sg_support}}

\begin{proof}[Proof of Lemma~\ref{lem:sg_support}]
For a fixed vector $v\in\mathbb R^{|\mathcal L|}$, consider the linear relaxation
\[
\max\left\{v^\top\delta:\ \mathbf 1^\top\delta\le K,\ 0\le \delta\le \mathbf 1\right\}.
\]
Its constraint matrix is totally unimodular, so every extreme point of the relaxation is binary.
Hence the relaxation is exact and has the same optimal value as
$\max_{\delta\in\Omega(K)} v^\top\delta$.

The dual of the relaxation is
\[
\min_{s\ge 0,\ u\ge 0}
\left\{Ks+\mathbf 1^\top u:\ u_\ell\ge v_\ell-s,\ \forall \ell\in\mathcal L\right\}.
\]
Strong duality then gives the desired identity.
\end{proof}

\subsection{Exactness of the linearized separation MILP}

The linearization in \eqref{eq:sg_sep_milp_main} is exact.
This fact is used by Algorithm~\ref{alg:sg_ccg} and by the convergence proof below.

\begin{proposition}[Exact linearized separation]
\label{prop:sg_sep_exact_appendix}
For any fixed $(x^{\star},\alpha^{\star},\lambda^{\star})$, the optimal value of the MILP
\eqref{eq:sg_sep_milp_main} is equal to
\[
\begin{aligned}
    &\mathrm{Viol}(x^{\star},\alpha^{\star},\lambda^{\star})=\\
    &
\max_{\boldsymbol\pi\in\Pi^B}
\left\{
\beta(\boldsymbol\pi)-\gamma(\boldsymbol\pi)^\top x^{\star}
+\max_{\delta\in\Omega(K)}\big(\phi(\boldsymbol\pi)-\lambda^{\star}\big)^\top\delta
-\alpha^{\star}
\right\}.
\end{aligned}
\]
\end{proposition}

\begin{proof}
Fix any tuple $\boldsymbol\pi=(\pi_1,\ldots,\pi_B)\in\Pi^B$ and define
\[
\omega:=\frac{1}{B}\sum_{b=1}^B U^\top\pi_b.
\]
By definition of the bounds in \eqref{eq:sg_bounds_main}, each component satisfies
$\underline\omega_\ell\le \omega_\ell\le \overline\omega_\ell$.
Now fix $\delta\in\Omega(K)$ and set $\tau_\ell:=\delta_\ell\omega_\ell$ for all
$\ell\in\mathcal L$.
Because $\delta_\ell\in\{0,1\}$ and $\omega_\ell\in[\underline\omega_\ell,\overline\omega_\ell]$,
the four McCormick inequalities in \eqref{eq:sg_sep_milp_main} are satisfied exactly.
Thus every feasible pair $(\boldsymbol\pi,\delta)$ of the original separation problem can be lifted to a feasible point
of \eqref{eq:sg_sep_milp_main} with the same objective value.
Consequently, the MILP optimal value is at least $\mathrm{Viol}(x^{\star},\alpha^{\star},\lambda^{\star})$.

Conversely, let $(\{\pi_b\},\delta,\omega,\tau)$ be any feasible solution of
\eqref{eq:sg_sep_milp_main}. Since $\delta_\ell$ is binary and the McCormick system is exact over a bounded interval,
those constraints imply
\[
\tau_\ell=\delta_\ell\omega_\ell,
\qquad \forall \ell\in\mathcal L.
\]
Using the defining equation for $\omega$ gives
\[
\mathbf 1^\top\tau
=\sum_{\ell\in\mathcal L}\delta_\ell\omega_\ell
=\delta^\top\left(\frac{1}{B}\sum_{b=1}^B U^\top\pi_b\right)
=\phi(\boldsymbol\pi)^\top\delta.
\]
Therefore the MILP objective becomes
\[
\frac{1}{B}\sum_{b=1}^B \pi_b^\top(H^b-Z^b x^{\star})
+\phi(\boldsymbol\pi)^\top\delta
-(\lambda^{\star})^\top\delta-\alpha^{\star},
\]
which is exactly the objective of the original separation problem.
Hence the MILP optimal value is at most $\mathrm{Viol}(x^{\star},\alpha^{\star},\lambda^{\star})$.
The two values are equal.
\end{proof}

\subsection{Finite exact reformulation}

The next observation is standard and is the only ingredient needed to prove finite termination.

\begin{proposition}[Finite exact reformulation]
\label{prop:sg_finite_exact_appendix}
Let $V:=\mathrm{ext}(\Pi)$ and, for each tuple
$r=(v_1,\ldots,v_B)\in V^B$, let
\[
\beta_r=\tfrac{1}{B}\sum_{b=1}^B v_b^\top H^b, \ 
\gamma_:=\tfrac{1}{B}\sum_{b=1}^B (Z^b)^\top v_b, \ 
\phi_r=\tfrac{1}{B}\sum_{b=1}^B U^\top v_b.
\]
Then \eqref{eq:sg_semiinf} is equivalent to the finite optimization problem
\begin{equation}
\label{eq:sg_finite_exact_appendix}
\begin{aligned}
\min \ 
&F^{\mathrm{cons}}(x)
+\tfrac{1-\pi_f}{A}\sum_{a=1}^A F^{{\rm nor},a}(x,y^{{\rm nor},a})
+\pi_f\big(\alpha+\lambda^\top \mathbf{FP}\big)\\
\text{s.t.}\ 
&x\in \mathcal{X},\qquad y^{{\rm nor},a}\in \mathcal{Y}^{\mathrm{nor}}(x,\xi^a),\quad a=1,\ldots,A,\\
&\lambda\ge 0,\\
&\alpha\ge \beta_r-\gamma_r^\top x+Ks_r+\mathbf 1^\top u_r,
\qquad \forall r\in V^B,\\
&u_{r,\ell}\ge \phi_{r,\ell}-\lambda_\ell-s_r,
\qquad \forall r\in V^B,\ \forall \ell\in\mathcal L,\\
&s_r\ge 0,\qquad u_r\ge 0,
\qquad \forall r\in V^B.
\end{aligned}
\end{equation}
\end{proposition}

\begin{proof}
For each sample $b$, the dual representation \eqref{eq:disaster_dual_compact} maximizes a linear function over the
polyhedron $\Pi$. Hence an optimal solution can always be chosen at an extreme point of $\Pi$.
Therefore the maximization in \eqref{eq:sg_fx_dual} may be restricted to tuples in $V^B$:
\[
f_{x,B}(\delta)
=
\max_{r\in V^B}
\left\{\beta_r-\gamma_r^\top x+\phi_r^\top\delta\right\}.
\]
Substituting this representation into Proposition~\ref{prop:sg_moment_dual} and applying
Lemma~\ref{lem:sg_support} to each tuple $r\in V^B$ gives \eqref{eq:sg_finite_exact_appendix}.
Since every step is exact, the two models are equivalent.
\end{proof}

\subsection{Proof of Theorem~\ref{thm:sg_finite}}

\begin{proof}[Proof of Theorem~\ref{thm:sg_finite}]
Let $z_B^{\star}$ denote the optimal value of the sampled problem \eqref{eq:overall_saa_compact},
or equivalently of the finite reformulation \eqref{eq:sg_finite_exact_appendix}. Let
$\mathcal R_k\subseteq V^B$ be the cut set used by the restricted master problem at iteration $k$,
and let $z_k^{\mathrm{RMP}}$ be its optimal value.
Because the RMP only contains a subset of the full cut family, it is a relaxation of
\eqref{eq:sg_finite_exact_appendix}; therefore,
\[
z_k^{\mathrm{RMP}}\le z_B^{\star}
\qquad \forall k.
\]

Now consider iteration $k$ and let $(x^k,\alpha^k,\lambda^k)$ be the RMP solution.
If the separation MILP \eqref{eq:sg_sep_milp_main} returns an optimal value at most
$\varepsilon_{\mathrm{cert}}=0$, then by Proposition~\ref{prop:sg_sep_exact_appendix} we have
\[
\mathrm{Viol}(x^k,\alpha^k,\lambda^k)\le 0.
\]
Hence no cut from the full family is violated, which means $(x^k,\alpha^k,\lambda^k)$ can be extended with suitable
support variables $(s_r,u_r)_{r\in V^B}$ to a feasible solution of the full model
\eqref{eq:sg_finite_exact_appendix}. Since the RMP is a relaxation of the full model, this implies
$z_k^{\mathrm{RMP}}=z_B^{\star}$, and the algorithm terminates optimally.

Suppose instead that the separation MILP returns a strictly positive optimal value. Let $\delta^k$ be an optimal
outage vector, and let $\widehat\pi_b^k\in V$ be the extreme optimal solution obtained by re-solving
\[
\max_{\pi_b\in\Pi}\ \pi_b^\top(H^b-Z^b x^k+U\delta^k),
\qquad b=1,\ldots,B.
\]
Set $r^k:=(\widehat\pi_1^k,\ldots,\widehat\pi_B^k)\in V^B$.
By construction, the cut associated with $r^k$ is violated at $(x^k,\alpha^k,\lambda^k)$.
Moreover, $r^k\notin\mathcal R_k$: if it were already present in the RMP, then the existing support variables
$(s_{r^k},u_{r^k})$ would make that cut satisfied, contradicting the positive violation.
Therefore each nonterminating iteration adds a previously unseen tuple from the finite set $V^B$.

Since $V$ is finite, so is $V^B$. Thus only finitely many distinct cuts can be added. The algorithm must therefore stop
after finitely many iterations, and by the first part of the proof the terminating solution is optimal for
\eqref{eq:overall_saa_compact}.
\end{proof}

\subsection{Proof of Theorem~\ref{prop:sg_epsopt}}

\begin{proof}[Proof of Theorem~\ref{prop:sg_epsopt}]
Let $(x^{\star},\alpha^{\star},\lambda^{\star})$ be the current RMP solution with objective value
$z^{\mathrm{RMP}}$. Assume that the separation routine provides a valid upper bound satisfying
\eqref{eq:sg_epscert}. Then for every $r\in V^B$,
\[
\beta_r-\gamma_r^\top x^{\star}
+\max_{\delta\in\Omega(K)}(\phi_r-\lambda^{\star})^\top\delta
-\alpha^{\star}
\le \varepsilon_{\mathrm{cert}}.
\]
Equivalently,
\[
\alpha^{\star}+\varepsilon_{\mathrm{cert}}
\ge
\beta_r-\gamma_r^\top x^{\star}
+\max_{\delta\in\Omega(K)}(\phi_r-\lambda^{\star})^\top\delta,
\qquad \forall r\in V^B.
\]
Now apply Lemma~\ref{lem:sg_support} to each vector $\phi_r-\lambda^{\star}$.
For every $r\in V^B$, there exist $\tilde s_r\ge 0$ and $\tilde u_r\ge 0$ such that
\[
\begin{aligned}
    \max_{\delta\in\Omega(K)}(\phi_r-\lambda^{\star})^\top\delta
=K\tilde s_r+\mathbf 1^\top \tilde u_r, \\
\tilde u_{r,\ell}\ge \phi_{r,\ell}-\lambda_\ell^{\star}-\tilde s_r,
\qquad \forall \ell\in\mathcal L.
\end{aligned}
\]
Hence
\[
(x^{\star},\alpha^{\star}+\varepsilon_{\mathrm{cert}},\lambda^{\star},\{\tilde s_r,\tilde u_r\}_{r\in V^B})
\]
is feasible for the full finite reformulation \eqref{eq:sg_finite_exact_appendix}. Its objective value is exactly
\[
z^{\mathrm{RMP}}+\pi_f\varepsilon_{\mathrm{cert}}.
\]
Therefore,
\[
z_B^{\star}\le z^{\mathrm{RMP}}+\pi_f\varepsilon_{\mathrm{cert}}.
\]
On the other hand, since the RMP is a relaxation of the full problem,
$z^{\mathrm{RMP}}\le z_B^{\star}$. Combining the two inequalities yields
\eqref{eq:sg_epsgap}.
\end{proof}

